\input{../tex-inputs/latex-leseansicht-vorspann}

\section[Arthur Schnitzler an Gustav Schwarzkopf, 12. 8. 1919]{L04526 Arthur Schnitzler an Gustav Schwarzkopf, 12. 8. 1919}
\nopagebreak\mylabel{L04526v}
\rehead{ }\normalsize\beginnumbering\briefempfaengerindex{Schwarzkopf, Gustav@\textsc{Schwarzkopf, Gustav}!zzzSchnitzler, Arthur@\emph{von Arthur Schnitzler}!1919-08-121@{12. 8. 1919}|(be}
\toendnotes[C]{\smallbreak\pagebreak[2]}
\correspDesc{Versand  durch Arthur Schnitzler am 12. 8. 1919 in Reichenau an der Rax
\newline{}Erhalt  durch Gustav Schwarzkopf im Zeitraum [13. 8. 1919 – 17. 8. 1919?] in Wien}\toendnotes[C]{\smallbreak}
\Standort{DLA, A:Schnitzler, HS.1985.1.1897.}
\physDesc{Bildpostkarte, 256 Zeichen
\newline{}Handschrift: Bleistift, lateinische Kurrent
\newline{}Versand: Stempel: »\nobreak{}\oindex{Reichenau an der Rax@\textbf{Reichenau an der Rax}|pwk}Reichenau, 12. VIII. \textcolor{gray}{19}, 5\nobreak{}«.  }\toendnotes[C]{\smallbreak}
\pstart
           \noindent{}\centering{}{\pb}\textcolor{gray}{\textbf{Reichenau (N.-Ö.)\oindex{Reichenau an der Rax@\textbf{Reichenau an der Rax}|pw} mit der Enge\oindex{Eng@\textbf{Eng}|pw}.}}\pend
           \pstart{}{\pb}Hrn Gustav Schwarzkopf\pend{}\pstart{}Wien I\oindex{I., Innere Stadt@\textbf{I., Innere Stadt}|pw}\pend{}\pstart{}Tiefer Graben 17\oindex{Wien@\textbf{Wien}!I., Innere Stadt@\textbf{I., Innere Stadt}!Tiefer Graben 17@\textbf{Tiefer Graben 17}|pw}\pend{}{\bigskip}\vspace{1em}
\pstart
           \raggedleft{}12. 8. 19\pend
           \vspace{0.5em}
\pstart
           Herzliche Grüße! Es ist – immerhin – sehr schön \introOben{}hier\introOben{}, und so
               verschieb ich  die Heimkehr auf \label{K_L04526-1v}\edtext{Ende der Woche}{\lemma{\textnormal{\emph{Ende der Woche}}}\Cendnote{\textnormal{Am 17. 8. 1919 kehrte Schnitzler 10 Tagen in Reichenau an der Rax\oindex{Reichenau an der Rax@\textbf{Reichenau an der Rax}|pwk} nach Wien\oindex{Wien@\textbf{Wien}|pwk} zurück.}}}\label{K_L04526-1}. Wir alle befinden uns wohl und beko{\geminationm}en ganz leidlich zu essen.\pend
           
\pstart
           Auf baldiges Wiedersehen!{\\[\baselineskip]}Ihr \spacefill\mbox{Arthur}\pend
           \leftskip=0em{}\selectlanguage{ngerman}\endnumbering\briefempfaengerindex{Schwarzkopf, Gustav@\textsc{Schwarzkopf, Gustav}!zzzSchnitzler, Arthur@\emph{von Arthur Schnitzler}!1919-08-121@{12. 8. 1919}|)be}\mylabel{L04526h}
\begin{anhang}
\end{anhang}\newcommand{\dateiname}{L04526}\newcommand{\titel}{Arthur Schnitzler an Gustav Schwarzkopf, 12. 8. 1919}\newcommand{\editorInnen}{Herausgegeben von Jahnke, SelmaMüller, Martin Anton}\input{../tex-inputs/latex-leseansicht-abspann}
